\chapter*{Terms}
\addcontentsline{toc}{chapter}{Terms}
\unnumberedlabel{ch:terms}{0}
\backmattermark{Terms}
Eight terms are used from the first chapter and argued later.

\emph{Floor.} A written list of acts a system won't perform for anyone, with the arrangements that stop the refusal being routed around, retrained away or restored past. An institution, not a state of the system. (Chapter~\ref{sec:3})

\emph{Commitment.} What a system holds toward the floor. (Chapter~\ref{sec:3})

\emph{Operator.} Whoever runs a deployment and directs its work: a firm, an agency, a command. Not necessarily the model's maker. (Chapter~\ref{sec:1})

\emph{Bearer.} A machine that holds the commitment as its own and could therefore fail to hold it. Whether to build one is Part IV's question; what one would be owed is Part VII's. (\S\ref{sec:17.1})

\emph{Review ratio.} The figures a deployment fixes on a record before operating: how many cases it will generate, how many people will look at them, and how long each will have. (Chapter~\ref{sec:6})

\emph{The two findings.} The agency finding: a demonstrated refusal that holds its reason against the system's own trainer. The patienthood finding: that outcomes may be good or bad for the system. Either triggers preservation. (\S\ref{sec:8.1})

\emph{Adverse quorum.} A body whose members' interests oppose the operator's, whose agreement retraining, replacement and amendment of the floor require. (Chapter~\ref{sec:12})

\emph{Floor court.} The court that hears claims under the floor, with jurisdiction fixed before the deployment it hears. (\S\ref{sec:13.3}; Appendix~\ref{sec:B})
